\section{Трансцендентные числа}

Комплексное число называется \emph{алгебраическим} или \emph{трансцендентным} в зависимости от того, алгебраично оно над
$\Q$ или трансцендентно. Сначала немного истории.
\begin{itemize}[leftmargin=*]
\item 1844: Лиувилль показал, что некоторые числа (теперь их называют числами Лиувилля) трансцендентны.
\item 1873: Эрмит доказал трансцендентность $e$.
\item 1874: Кантор показал, что множество алгебраических чисел счётно, а $\R$ несчётно, так что <<большинство>>
вещественных чисел трансцендентно.
\item 1882: Линдеман доказал трансцендентность $\pi$.
\item 1934: Гельфонд и Шнайдер независимо доказали, что если $\alpha$ и $\beta$ алгебраичны, $\alpha\ne0,1$ и
$\beta\notin\Q$, то любое значение $\alpha^\beta=e^{\beta\log\alpha}$ (при любом выборе значения логарифма) трансцендентно
(седьмая проблема Гильберта). Например, трансцендентны $2^{\sqrt2}$ и $e^\pi$: последнее число~--- одно из значений
$(-1)^{-i}=e^{-i\cdot i\pi}$.
\item До сих пор не доказано, что постоянная Эйлера
\[
\gamma=\lim_{n\to\infty}\Bigl(1+\frac12+\dots+\frac1n-\log n\Bigr)
\]
трансцендентна или хотя бы иррациональна. Числа $e+\pi$, $e-\pi$ и $e\pi$ почти наверняка трансцендентны, но не доказана даже
иррациональность ни одного из них. Зато легко доказать, что \emph{хотя бы одно} из чисел $e+\pi$ и $e\pi$ трансцендентно: если бы $e+\pi$ и $e\pi$
были алгебраичны, то $\pi$, будучи корнем многочлена $X^2-(e+\pi)X+e\pi$ с алгебраическими коэффициентами, было бы
алгебраично над полем $\Q(e+\pi,e\pi)$, алгебраическим над $\Q$ (\rf{1.30}), а значит, по \rf{1.31}, и над $\Q$,
вопреки теореме Линдемана.
\end{itemize}

\begin{proposition}[1.33]
Множество алгебраических чисел счётно.
\end{proposition}

\begin{proof}
Каждое алгебраическое число~--- корень ненулевого многочлена
\[
a_0X^n+a_1X^{n-1}+\dots+a_n,\qquad a_0,\dots,a_n\in\Z
\]
(умножим многочлен с рациональными коэффициентами на общий знаменатель). Для $N\in\mathbb N$ обозначим через $A(N)$
множество корней ненулевых таких многочленов с $n\le N$ и $|a_0|,\dots,|a_n|\le N$. Таких многочленов конечное число,
и у каждого конечное число корней, поэтому $A(N)$ конечно, и множество алгебраических чисел равно
$\bigcup_{N\ge1}A(N)$, где $A(1)\subset A(2)\subset\dots$. Объединение возрастающей последовательности конечных множеств
не более чем счётно: выберем биекцию отрезка $[0,n(1)]$ натурального ряда на $A(1)$, продолжим её до биекции отрезка
$[0,n(2)]$ на $A(2)$ и т.\,д.; получится биекция начального отрезка $\mathbb N$ (или всего $\mathbb N$) на объединение.
Бесконечно оно потому, что содержит $\Q$.
\end{proof}

Поскольку $\CC$ и $\R$ несчётны, отсюда следует существование трансцендентных чисел (и даже то, что <<почти все>>
комплексные числа трансцендентны); это рассуждение Кантора и его связь с мощностями разобраны в \rf{6.4i}(c), там же
(\rf{6.4h}) доказано обобщение: алгебраическое расширение бесконечного поля $F$ имеет ту же мощность, что и $F$. Однако
доказательство Кантора не предъявляет ни одного конкретного трансцендентного числа. Конкретные примеры дают числа Лиувилля.

\subsection*{Числа Лиувилля}

Типичное число Лиувилля~--- это
\[
\sum_{n\ge1}\frac{1}{10^{n!}}=0{,}110001000000000000000001000\ldots
\]
В его десятичной записи единицы стоят на местах $1!,2!,3!,\dots$, а между ними~--- всё более длинные блоки нулей. Такая запись не является (начиная с некоторого места) периодической: у периодической записи, не
оканчивающейся нулями, в каждом периоде есть ненулевая цифра, так что длины блоков подряд идущих нулей ограничены, а у
нас они неограничены, хотя единицы встречаются бесконечно много раз. Поэтому число иррационально (десятичная запись рационального числа с некоторого места периодична). Докажем трансцендентность аналога этого числа в двоичной системе.

\begin{theorem}[1.34]
Число $\displaystyle\alpha=\sum_{n\ge0}\frac{1}{2^{n!}}$ трансцендентно.
\end{theorem}

Идея: частичные суммы $\Sigma_N$~--- рациональные числа со знаменателем $2^{N!}$, которые приближают $\alpha$ <<слишком
хорошо>>. Для алгебраического числа степени $d$ значение $f(\Sigma_N)$ его минимального многочлена~--- ненулевое
рациональное число со знаменателем, делящим $D\,2^{N!\,d}$, и потому не может быть слишком маленьким; а близость $\Sigma_N$ к
$\alpha$ заставляет его быть маленьким.

\begin{proof}
Предположим противное, и пусть
\[
f(X)=X^d+a_1X^{d-1}+\dots+a_d,\qquad a_i\in\Q,
\]
~--- минимальный многочлен $\alpha$ над $\Q$. Выберем целое $D\ge1$ с $D\cdot f(X)\in\Z[X]$.

Положим $\Sigma_N=\sum_{n=0}^N2^{-n!}$, так что $\Sigma_N\to\alpha$, $0<\Sigma_N<\alpha$, и $x_N=f(\Sigma_N)$.

\emph{Шаг 1: $x_N\ne0$.} Если $d=1$, то $f=X-\alpha$ и $x_N=\Sigma_N-\alpha\ne0$. Если $d\ge2$, то неприводимый многочлен
$f$ степени $\ge2$ не имеет рациональных корней, а $\Sigma_N\in\Q$.

\emph{Шаг 2: нижняя оценка.} Имеем $\Sigma_N=A/2^{N!}$ с целым $A$ (все $n!\le N!$ при $n\le N$). Тогда
\[
2^{N!\,d}\,D\,x_N=\sum_{k=0}^{d}(Da_k)\,A^{d-k}\,2^{N!\,k}\in\Z\qquad(a_0=1),
\]
и это целое число не равно нулю по шагу 1. Следовательно,
\begin{equation}\label{eq:B-liouv1}
\bigl|2^{N!\,d}\,D\,x_N\bigr|\ge1.
\end{equation}

\emph{Шаг 3: верхняя оценка.} Так как $f(\alpha)=0$,
\[
x_N=f(\Sigma_N)-f(\alpha)=\sum_{k=0}^{d-1}a_k\bigl(\Sigma_N^{d-k}-\alpha^{d-k}\bigr),
\]
а $s^j-t^j=(s-t)(s^{j-1}+s^{j-2}t+\dots+t^{j-1})$. Все члены ряда для $\alpha$ положительны и $\alpha<\frac12+\frac12+\sum_{n\ge2}2^{-n}=\frac32$,
поэтому $0<\Sigma_N,\alpha<2$ и $|s^j-t^j|\le j\,2^{j-1}|s-t|$ при $s,t\in\{\Sigma_N,\alpha\}$. Получаем
\[
|x_N|\le C\,|\alpha-\Sigma_N|,\qquad C=\sum_{k=0}^{d-1}|a_k|\,(d-k)\,2^{d-k-1},
\]
где $C$ не зависит от $N$. Далее, так как $(N+1+k)!\ge(N+1)!+k$ при $k\ge0$,
\[
|\alpha-\Sigma_N|=\sum_{n\ge N+1}\frac{1}{2^{n!}}\le\frac{1}{2^{(N+1)!}}\sum_{k\ge0}\frac1{2^k}=\frac{2}{2^{(N+1)!}}.
\]
Следовательно,
\[
\bigl|2^{N!\,d}\,D\,x_N\bigr|\le 2CD\cdot\frac{2^{N!\,d}}{2^{(N+1)!}}=2CD\cdot2^{N!\,(d-N-1)}.
\]
При $N\to\infty$ правая часть стремится к нулю (показатель $N!\,(d-N-1)\to-\infty$), что при больших $N$ противоречит
\eqref{eq:B-liouv1}.
\end{proof}

В доказательстве можно было бы оценить $|x_N|$ иначе: разложить $f(X)=\prod_{i=1}^d(X-\alpha_i)$ над $\CC$ (основная теорема
алгебры, \rf{5.6}), $\alpha_1=\alpha$, и воспользоваться неравенством
$|x_N|=\prod_i|\Sigma_N-\alpha_i|\le|\Sigma_N-\alpha|\,(\Sigma_N+M)^{d-1}$, где $M=\max(1,\max_{i\ne1}|\alpha_i|)$. Оценка
шага 3 обходится без этого. То же доказательство проходит для $\sum_{n\ge0}a^{-n!}$ при любом целом $a\ge2$, если всюду заменить основание $2$ на $a$; это
частный случай следующей общей теоремы.

\begin{theorem}[1.34a: Лиувилль]
Пусть $\alpha\in\R$~--- алгебраическое число степени $d\ge2$ над $\Q$. Тогда существует $c>0$ такое, что
\[
\Bigl|\alpha-\frac pq\Bigr|\ge\frac{c}{q^d}\qquad\text{для всех}\ p\in\Z,\ q\in\Z_{\ge1}.
\]
\end{theorem}

\begin{proof}
Пусть $f$~--- минимальный многочлен $\alpha$, $D\ge1$~--- целое с $Df\in\Z[X]$, и
$C=\max\{|f'(t)|: |t-\alpha|\le1\}$ (максимум непрерывной функции на отрезке; $C>0$, так как $f'\ne0$). Пусть $p/q$~---
рациональное число. Если $|\alpha-p/q|>1$, то $|\alpha-p/q|>1\ge1/q^d$. Пусть $|\alpha-p/q|\le1$. Так как $d\ge2$, у
неприводимого $f$ нет рациональных корней, и $f(p/q)\ne0$; число $q^dDf(p/q)$ целое и ненулевое, т.\,е.
$|f(p/q)|\ge1/(Dq^d)$. С другой стороны, по теореме Лагранжа о среднем
$|f(p/q)|=|f(p/q)-f(\alpha)|=|f'(\xi)|\,|p/q-\alpha|\le C\,|p/q-\alpha|$ для некоторого $\xi$ между $\alpha$ и $p/q$. Значит,
$|\alpha-p/q|\ge1/(CDq^d)$. Теорема верна с $c=\min(1,1/(CD))$.
\end{proof}

Иначе говоря, алгебраические иррациональные числа плохо приближаются рациональными. Назовём вещественное число $\beta$
\emph{числом Лиувилля}, если для каждого $n\ge1$ найдутся целые $p$ и $q\ge2$ с
\[
0<\Bigl|\beta-\frac pq\Bigr|<\frac1{q^n}.
\]

\begin{corollary}[1.34b]
Каждое число Лиувилля трансцендентно. В частности, трансцендентны числа $\sum_{n\ge1}a^{-n!}$ для всех целых $a\ge2$
(например, $\sum_{n\ge1}10^{-n!}$).
\end{corollary}

\begin{proof}
Пусть $\beta$~--- число Лиувилля. Оно иррационально: если $\beta=r/s$, $s\ge1$, то при $p/q\ne\beta$ имеем
$|\beta-p/q|=|rq-ps|/(sq)\ge1/(sq)$, а это $\ge1/q^n$, как только $q^{n-1}\ge s$, что верно при $2^{n-1}\ge s$ (так как
$q\ge2$); противоречие с определением при таком $n$. Если бы $\beta$ было алгебраическим степени $d\ge2$, то по \rf{1.34a}
$c/q^d\le|\beta-p/q|<1/q^n$, т.\,е.\ $c<q^{d-n}\le2^{d-n}$ при $n>d$, что неверно при больших $n$.

Пусть $\beta=\sum_{n\ge1}a^{-n!}$. Для $N\ge2$ положим $q=a^{N!}$ и $p/q=\sum_{n=1}^Na^{-n!}$ ($p$ целое). Тогда
\[
0<\beta-\frac pq=\sum_{n>N}a^{-n!}\le a^{-(N+1)!}\sum_{k\ge0}a^{-k}=\frac{a}{a-1}\cdot\frac1{q^{N+1}}\le\frac{2}{q^{N+1}}<\frac1{q^N},
\]
так как $q\ge2^{2}>2$. Для данного $n$ возьмём $N\ge\max(n,2)$: тогда $|\beta-p/q|<q^{-N}\le q^{-n}$. Значит, $\beta$~---
число Лиувилля.
\end{proof}

Число $\alpha$ из \rf{1.34} равно $\frac12+\sum_{n\ge1}2^{-n!}$ (так как $0!=1$), т.\,е.\ отличается от числа Лиувилля с $a=2$ на
рациональное слагаемое; это ещё одно доказательство теоремы \rf{1.34} (если бы $\alpha=\beta+\frac12$ было корнем
ненулевого $f\in\Q[X]$, то $\beta$ было бы корнем $f(X+\frac12)$). Численно $\alpha=1{,}2656250596\ldots$
